\item \textbf{Problema de repaso.} Suponga que cierto líquido, con densidad de $1230 \text{ kg/m}^3$, no ejerce fuerza de fricción sobre objetos esféricos. Una bola de $2.10 \text{ kg}$ de masa y $9.00 \text{ cm}$ de radio se deja caer desde el reposo en un tanque profundo de este líquido desde una altura de $3.30 \text{ m}$ sobre la superficie. (a) Encuentre la rapidez con que entra la bola al líquido. (b) Evalúe las magnitudes de las dos fuerzas que se ejercen sobre la bola mientras se mueve a través del líquido. (c) Explique por qué la bola se mueve hacia abajo sólo una distancia limitada en el líquido y calcule esta distancia. (d) ¿Con qué rapidez la bola aparece fuera del líquido? (e) ¿De qué modo se compara el intervalo de tiempo $\Delta t_{\text{abajo}}$, durante el cual la bola se mueve desde la superficie hasta su punto más bajo, con el intervalo de tiempo $\Delta t_{\text{arriba}}$ para el viaje de regreso entre los mismos puntos? (f) ¿Qué pasaría si? Ahora modifique el modelo para suponer que el líquido ejerce una pequeña fuerza de fricción sobre la bola, opuesta en dirección a su movimiento. En este caso, ¿de qué forma se comparan los intervalos de tiempo $\Delta t_{\text{abajo}}$ y $\Delta t_{\text{arriba}}$? Explique su respuesta con un argumento conceptual en lugar de un cálculo numérico.
